% Source: ME33_v1.html
\documentclass[11pt]{article}
\usepackage[utf8]{inputenc}
\usepackage[T1]{fontenc}
\usepackage[margin=1in]{geometry}
\usepackage{amsmath}
\usepackage{amssymb}
\usepackage{siunitx}
\usepackage{enumitem}
\usepackage{parskip}

\title{PS160 -- ME33\_v1.html}
\date{}

\begin{document}
\maketitle

\section*{Questions}

\begin{enumerate}[leftmargin=*]
  \item The figure (not to scale) shows two plane mirrors and an incident light ray. Find the angle $\alpha$ if $\phi$ = 62$^{\circ}$ and $\theta$ = 28$^{\circ}$.

[Image: ME33\_plane\_mirror\_2.png]

$\alpha$ = $^{\circ}$ \hfill \textit{(10 pts, Numeric)}

  \item Suppose the index of refraction of a certain medium is 1.91. Calculate the speed of light in that medium and supply the numerical prefix missing from the blank below.

\emph{v} = $\times$ 10\textsuperscript{8} m/s \hfill \textit{(10 pts, Numeric)}

  \item Light of frequency 6.23 $\times$ 10\textsuperscript{14} Hz passes from air into a material having an index of refraction 1.21. What is the wavelength (in nm) of the light in this material? \hfill \textit{(10 pts, Numeric)}

  \item A ray of light traveling in air is incident at an unknown angle with respect to the normal on a fluid with index of refraction 1.61. If the refraction angle is 21$^{\circ}$, calculate the angle of incidence in degrees. \hfill \textit{(10 pts, Numeric)}

  \item The figure shows a light ray incident on a layer of oil that is floating on water. Calculate the angle $\phi$ in degrees if \emph{n}\textsubscript{oil} = 1.53, \emph{n}\textsubscript{water} = 1.33, and $\theta$ = 32$^{\circ}$.

[Image: three\_layers\_01.png] \hfill \textit{(10 pts, Numeric)}

  \item The table below shows the index of refraction for BK7 optical glass at several different wavelengths. White light is incident on the top face of a large block of this glass at an angle of 48$^{\circ}$ with respect to the normal. If the block is 0.42 m thick, by how many millimeters are beams at wavelengths 2 and 4 separated at the bottom face of the block?

Number
Wavelength (nm)
Index of Refraction

1
400
1.531

2
475
1.523

3
550
1.519

4
625
1.515 \hfill \textit{(10 pts, Numeric)}

  \item A ray of light is traveling in a substance with index of refraction \emph{n}\textsubscript{1} when it encounters a boundary with a second material having index of refraction \emph{n}\textsubscript{2} = 1.5, as in the figure where $\theta$ = 49$^{\circ}$. Calculate \emph{n}\textsubscript{1}.

[Image: total\_internal\_reflection\_1.png] \hfill \textit{(10 pts, Numeric)}

  \item A liquid-air boundary has a critical angle of 42.3$^{\circ}$ for total internal reflection. If a ray of light traveling in air has an incident angle of 34$^{\circ}$ at the liquid-air boundary, calculate the angle between the refracted ray and the surface normal. \hfill \textit{(10 pts, Numeric)}

  \item Monochromatic, polarized light passes through a polarization filter with a polarization axis that is rotated 37$^{\circ}$ with respect to the light polarization. The light then passes through a second polarization filter with axis rotated a further 49$^{\circ}$ from that of the previous filter. Calculate the ratio of the final light intensity to the original light intensity. \hfill \textit{(10 pts, Numeric)}

  \item A certain pane of glass has a refractive index of 1.6. For what incident angle in degrees is reflected light completely polarized if the glass is immersed in medium number 3 from the table below?

Number
Medium
Index of Refraction

1
air
1.000

2
water
1.333

3
oil
1.465 \hfill \textit{(10 pts, Numeric)}

  \item In general, the [DROPDOWN: smaller | larger]
the difference in index of refraction, the more an incident ray of light bends at a boundary. The index of refraction, however, depends on [DROPDOWN: frequency | the angle of incidence | wavelength]
, which is called [DROPDOWN: desiccation | dispersion | dissipation]
. Shorter wavelengths bend [DROPDOWN: more than | the same as | less than]
longer wavelengths. \hfill \textit{(10 pts, Dropdown)}

  \item Total internal reflection can only occur when a ray of light encounters a medium with a [DROPDOWN: higher | lower]
index of refraction. For this reason, a ray traveling in a medium with index of refraction [DROPDOWN: greater than one | less than one | equal to one]
cannot undergo total internal reflection when encountering the boundary of another medium. \hfill \textit{(10 pts, Dropdown)}

\end{enumerate}

\end{document}